% unit: u19
% sources: sources/089-orthogonal-matrices-and-gram-schmidt.txt, sources/032-lecture-17-orthogonal-matrices-and-gram-schmidt.txt, sources/182-mit18-06scf11-ses2-4sum-pdf-download-file.txt, sources/033-problem-solving-gram-schmidt-orthogonalization.txt, sources/183-mit18-06scf11-ses2-4prob-pdf-download-file.txt, sources/184-mit18-06scf11-ses2-4sol-pdf-download-file.txt
\section{Orthogonal matrices and Gram--Schmidt}\label{sec:u19}

The last step in the story of orthogonality is an orthogonal \emph{basis}. Section~\ref{sec:u18}
ended with the observation that the normal equations become trivial when the columns of $A$
are perpendicular unit vectors, because then $A\T A=I$. This section collects such matrices,
called $Q$, shows exactly which formulas simplify, and then gives the Gram--Schmidt process,
which turns any independent columns into orthonormal ones. In matrix language Gram--Schmidt is
a factorization $A=QR$, just as elimination was $A=LU$.

\subsection{Orthonormal vectors and the matrix \texorpdfstring{$Q$}{Q}}

\begin{definition}[Orthonormal vectors]\label{def:u19-orthonormal}
Vectors $\vb{q}_1,\dots,\vb{q}_n$ in $\R^m$ are \emph{orthonormal} if
\[
  \vb{q}_i\T\vb{q}_j=\begin{cases}0 & i\neq j\quad\text{(orthogonal)},\\ 1 & i=j\quad\text{(normalized: unit length)}.\end{cases}
\]
\end{definition}

Put orthonormal vectors into the columns of an $m\times n$ matrix $Q$. Then $Q\T Q$ is the
matrix of all their dot products: row $i$ of $Q\T$ is $\vb{q}_i\T$, column $j$ of $Q$ is
$\vb{q}_j$.

\begin{proposition}\label{prop:u19-qtq}
$Q$ has orthonormal columns if and only if $Q\T Q=I$ (the $n\times n$ identity).
\end{proposition}

\begin{proof}
The $(i,j)$ entry of $Q\T Q$ is $\vb{q}_i\T\vb{q}_j$; Definition~\ref{def:u19-orthonormal}
says precisely that these are the entries of $I$.
\end{proof}

$Q$ need not be square: two orthonormal columns in $\R^4$ give a $4\times 2$ matrix with
$Q\T Q=I_2$ (but $QQ\T\neq I_4$). Orthonormal columns are automatically independent
(Exercise~\ref{ex:u19-indep}). This adds a new family to our list of special matrices
(triangular, diagonal, permutation, symmetric, echelon, projection).

\begin{definition}[Orthogonal matrix]\label{def:u19-orthogonal}
A \emph{square} matrix with orthonormal columns is called an \emph{orthogonal matrix}. For
such $Q$, $Q\T Q=I$ means $Q\T=Q^{-1}$, and then also $QQ\T=I$.
\end{definition}

The name is a historical accident (``orthonormal matrix'' would be more accurate), and it is
reserved for the square case.

\begin{example}[A gallery of orthogonal matrices]\label{ex:u19-gallery}
Four families are worth knowing.
\begin{itemize}
  \item Every permutation matrix, e.g.\
    $Q=\left[\begin{smallmatrix}0&0&1\\1&0&0\\0&1&0\end{smallmatrix}\right]$: its columns
    are distinct standard basis vectors, and $Q\T$ (another permutation) is its inverse.
  \item Rotations
    $Q=\begin{bmatrix}\cos\theta&-\sin\theta\\ \sin\theta&\cos\theta\end{bmatrix}$: both
    columns have length $\cos^2\theta+\sin^2\theta=1$ and their dot product is
    $-\cos\theta\sin\theta+\sin\theta\cos\theta=0$.
  \item $\begin{bmatrix}1&1\\1&-1\end{bmatrix}$ has orthogonal columns of length $\sqrt2$;
    dividing by $\sqrt2$ makes it orthogonal. Repeating the sign pattern in blocks gives
    \[
      Q=\frac12\begin{bmatrix}1&1&1&1\\1&-1&1&-1\\1&1&-1&-1\\1&-1&-1&1\end{bmatrix},
    \]
    a \emph{Hadamard matrix}: entries $\pm1$ (scaled), columns of length $2$ before
    scaling. Such matrices exist for sizes $2,4,8,16,\dots$ and many others, but which sizes
    allow them is still an open question.
  \item A rectangular example without square roots, and a square one extending it:
    \[
      Q=\frac13\begin{bmatrix}1&-2\\2&-1\\2&2\end{bmatrix},\qquad
      Q'=\frac13\begin{bmatrix}1&-2&2\\2&-1&-2\\2&2&1\end{bmatrix}.
    \]
    Each column has squared length $\frac{1+4+4}{9}=1$, and the dot products vanish, e.g.\
    $(1,2,2)\cdot(-2,-1,2)=-2-2+4=0$.
\end{itemize}
\end{example}

\subsection{Orthonormal columns make projection easy}

Suppose $Q$ has orthonormal columns and we project onto its column space. The projection
formula of Section~\ref{sec:u17} contains $Q\T Q=I$, which disappears:

\begin{proposition}[Projection with an orthonormal basis]\label{prop:u19-proj}
If $Q$ has orthonormal columns, the projection onto $\Col{Q}$ is
\[
  P=Q(Q\T Q)^{-1}Q\T=QQ\T ,
\]
the least-squares solution of $Q\vb{x}=\vb{b}$ is $\hat{\vb{x}}=Q\T\vb{b}$, that is
$\hat{x}_i=\vb{q}_i\T\vb{b}$, and
$\vb{p}=\vb{q}_1(\vb{q}_1\T\vb{b})+\cdots+\vb{q}_n(\vb{q}_n\T\vb{b})$.
If $Q$ is square, $P=I$.
\end{proposition}

\begin{proof}
Substitute $Q\T Q=I$ into $P=A(A\T A)^{-1}A\T$ and into the normal equations
$Q\T Q\hat{\vb{x}}=Q\T\vb{b}$. Writing $QQ\T\vb{b}$ column by column gives the sum of
one-dimensional projections. If $Q$ is square, $QQ\T=I$ by Definition~\ref{def:u19-orthogonal}.
As a check, $QQ\T$ is symmetric and $(QQ\T)(QQ\T)=Q(Q\T Q)Q\T=QQ\T$.
\end{proof}

No matrix needs to be inverted and no system solved: each coefficient is one dot product.
This is the formula behind Fourier series and much of numerical linear algebra, where
orthonormal vectors are prized because they never grow or shrink, so roundoff stays under
control. For the $3\times 2$ matrix of Example~\ref{ex:u19-gallery},
\[
  QQ\T=\frac19\begin{bmatrix}5&4&-2\\4&5&2\\-2&2&8\end{bmatrix},
\]
the projection onto the plane spanned by $(1,2,2)$ and $(-2,-1,2)$.

\subsection{The Gram--Schmidt process}

Elimination aimed for a triangular matrix; Gram--Schmidt aims for orthonormal columns. Start
with two independent vectors $\vb{a}$ and $\vb{b}$. Keep the direction of the first:
$\vb{A}=\vb{a}$. The second is not perpendicular to $\vb{A}$, so remove its projection onto
$\vb{A}$ and keep the rest, which is exactly the error vector $\vb{e}$ of
Section~\ref{sec:u17}:
\[
  \vb{B}=\vb{b}-\frac{\vb{A}\T\vb{b}}{\vb{A}\T\vb{A}}\,\vb{A}.
\]
It is perpendicular to $\vb{A}$:
$\vb{A}\T\vb{B}=\vb{A}\T\vb{b}-\frac{\vb{A}\T\vb{b}}{\vb{A}\T\vb{A}}\vb{A}\T\vb{A}=0$
(Figure~\ref{fig:u19-gs}). Finally divide by lengths:
$\vb{q}_1=\vb{A}/\|\vb{A}\|$, $\vb{q}_2=\vb{B}/\|\vb{B}\|$. With a third vector $\vb{c}$,
subtract its components along \emph{both} earlier directions:
\[
  \vb{C}=\vb{c}-\frac{\vb{A}\T\vb{c}}{\vb{A}\T\vb{A}}\,\vb{A}-\frac{\vb{B}\T\vb{c}}{\vb{B}\T\vb{B}}\,\vb{B},
  \qquad \vb{q}_3=\vb{C}/\|\vb{C}\|.
\]

\begin{figure}[ht]
  \centering
  \begin{tikzpicture}[>=Stealth,scale=1.3]
    \draw[->,very thick,blue] (0,0) -- (3,0.8) node[right] {$\vb{a}=\vb{A}$};
    \draw[->,very thick,red!75!black] (0,0) -- (1.5,2.2) node[above] {$\vb{b}$};
    \draw[->,thick,green!50!black] (0,0) -- (1.948,0.5195) node[below right=-1pt] {$\frac{\vb{A}\T\vb{b}}{\vb{A}\T\vb{A}}\vb{A}$};
    \draw[dashed] (1.948,0.5195) -- (1.5,2.2);
    \draw[->,very thick,orange!85!black] (0,0) -- (-0.448,1.6805) node[above left] {$\vb{B}=\vb{b}-\frac{\vb{A}\T\vb{b}}{\vb{A}\T\vb{A}}\vb{A}$};
    \draw[rotate=14.93] (0.2,0) -- (0.2,0.2) -- (0,0.2);
    \draw[shift={(1.948,0.5195)},rotate=14.93] (-0.18,0) -- (-0.18,0.18) -- (0,0.18);
    \fill (0,0) circle (1.2pt) node[below left] {$\vb{0}$};
  \end{tikzpicture}
  \caption{One Gram--Schmidt step. Subtracting from $\vb{b}$ its projection onto $\vb{A}$
  leaves $\vb{B}$, perpendicular to $\vb{A}$ and parallel to the dashed error segment. The
  plane spanned by $\vb{A},\vb{B}$ is the plane spanned by $\vb{a},\vb{b}$.}
  \label{fig:u19-gs}
\end{figure}

\begin{theorem}[Gram--Schmidt]\label{thm:u19-gs}
Let $\vb{a}_1,\dots,\vb{a}_n$ be independent. Define recursively
\[
  \vb{A}_k=\vb{a}_k-\sum_{i<k}(\vb{q}_i\T\vb{a}_k)\,\vb{q}_i,\qquad
  \vb{q}_k=\vb{A}_k/\|\vb{A}_k\|.
\]
Then each $\vb{A}_k\neq\vb{0}$, the vectors $\vb{q}_1,\dots,\vb{q}_n$ are orthonormal, and for
every $k$
\[
  \Span\{\vb{q}_1,\dots,\vb{q}_k\}=\Span\{\vb{a}_1,\dots,\vb{a}_k\}.
\]
\end{theorem}

(Since $\vb{q}_i=\vb{A}_i/\|\vb{A}_i\|$, the term $(\vb{q}_i\T\vb{a}_k)\vb{q}_i$ equals
$\frac{\vb{A}_i\T\vb{a}_k}{\vb{A}_i\T\vb{A}_i}\vb{A}_i$, the form used above.)

\begin{proof}
Induction on $k$. Assume $\vb{q}_1,\dots,\vb{q}_{k-1}$ are orthonormal and span the same
space $V_{k-1}$ as $\vb{a}_1,\dots,\vb{a}_{k-1}$. The vector $\vb{A}_k$ is $\vb{a}_k$ minus an
element of $V_{k-1}$. If $\vb{A}_k$ were $\vb{0}$, then $\vb{a}_k\in V_{k-1}$, contradicting
independence; so $\vb{q}_k$ is defined. For $j<k$,
\[
  \vb{q}_j\T\vb{A}_k=\vb{q}_j\T\vb{a}_k-\sum_{i<k}(\vb{q}_i\T\vb{a}_k)\,\vb{q}_j\T\vb{q}_i
  =\vb{q}_j\T\vb{a}_k-\vb{q}_j\T\vb{a}_k=0,
\]
because only the $i=j$ term survives. So $\vb{q}_k$ is a unit vector orthogonal to the
earlier $\vb{q}$'s. Finally $\vb{q}_k$ is a combination of $\vb{a}_k$ and $V_{k-1}$, and
$\vb{a}_k=\|\vb{A}_k\|\vb{q}_k+\sum_{i<k}(\vb{q}_i\T\vb{a}_k)\vb{q}_i$ is a combination of
$\vb{q}_1,\dots,\vb{q}_k$; hence the two spans agree.
\end{proof}

\begin{example}[Two vectors in \texorpdfstring{$\R^3$}{R3}]\label{ex:u19-lecture}
Let $\vb{a}=(1,1,1)$ and $\vb{b}=(1,0,2)$. Then $\vb{A}=\vb{a}$, $\vb{A}\T\vb{b}=3$,
$\vb{A}\T\vb{A}=3$, and
\[
  \vb{B}=\begin{bmatrix}1\\0\\2\end{bmatrix}-\frac33\begin{bmatrix}1\\1\\1\end{bmatrix}
  =\begin{bmatrix}0\\-1\\1\end{bmatrix},
  \qquad \vb{A}\T\vb{B}=0-1+1=0 .
\]
Dividing by the lengths $\sqrt3$ and $\sqrt2$,
\[
  Q=\begin{bmatrix}\vb{q}_1&\vb{q}_2\end{bmatrix}
  =\begin{bmatrix}1/\sqrt3&0\\1/\sqrt3&-1/\sqrt2\\1/\sqrt3&1/\sqrt2\end{bmatrix}.
\]
The column space of $Q$ is the same plane as the column space of
$\left[\begin{smallmatrix}1&1\\1&0\\1&2\end{smallmatrix}\right]$; only the basis improved.
\end{example}

\begin{trap}
When computing $\vb{C}$, subtract the projections of $\vb{c}$ onto the \emph{new} orthogonal
vectors $\vb{A}$ and $\vb{B}$, not onto the original $\vb{a}$ and $\vb{b}$; the original
vectors are not perpendicular, and subtracting both of their projections over-corrects.
And do not forget the final division by lengths: $\vb{A},\vb{B},\vb{C}$ are orthogonal but
generally not orthonormal. Doing the normalization last keeps square roots out of the
intermediate arithmetic.
\end{trap}

\subsection{\texorpdfstring{$A=QR$}{A = QR}}

Elimination is summarized by $A=LU$; Gram--Schmidt is summarized by a factorization that
connects $A$ to $Q$. Since $\Col{A}=\Col{Q}$, each column of $A$ is a combination of the
$\vb{q}$'s, and the coefficients form a matrix $R$.

\begin{theorem}[\texorpdfstring{$QR$}{QR} factorization]\label{thm:u19-qr}
If $A$ ($m\times n$) has independent columns and $Q$ comes from Gram--Schmidt, then
\[
  A=QR,\qquad R=Q\T A,\qquad r_{ij}=\vb{q}_i\T\vb{a}_j ,
\]
and $R$ is upper triangular with positive diagonal entries $r_{kk}=\|\vb{A}_k\|$.
\end{theorem}

\begin{proof}
Rearranging the definition of $\vb{A}_k$,
$\vb{a}_k=\sum_{i<k}(\vb{q}_i\T\vb{a}_k)\vb{q}_i+\|\vb{A}_k\|\,\vb{q}_k$. This says column $k$
of $A$ equals $Q$ times the column $(r_{1k},\dots,r_{kk},0,\dots,0)$, with
$r_{ik}=\vb{q}_i\T\vb{a}_k$ for $i<k$ and $r_{kk}=\|\vb{A}_k\|>0$; note also
$\vb{q}_k\T\vb{a}_k=\|\vb{A}_k\|$ by dotting the same identity with $\vb{q}_k$. The entries
below the diagonal are zero because $\vb{a}_k$ lies in the span of
$\vb{q}_1,\dots,\vb{q}_k$, so $\vb{q}_i\T\vb{a}_k=0$ for $i>k$: later $\vb{q}$'s are
perpendicular to earlier $\vb{a}$'s. Multiplying $A=QR$ by $Q\T$ gives $R=Q\T A$.
\end{proof}

For two columns,
\[
  \begin{bmatrix}\vb{a}_1&\vb{a}_2\end{bmatrix}
  =\begin{bmatrix}\vb{q}_1&\vb{q}_2\end{bmatrix}
  \begin{bmatrix}\vb{q}_1\T\vb{a}_1&\vb{q}_1\T\vb{a}_2\\0&\vb{q}_2\T\vb{a}_2\end{bmatrix},
\]
and the zero is $\vb{q}_2\T\vb{a}_1=0$ because $\vb{q}_2$ was built perpendicular to
$\vb{a}_1=\vb{A}$. In Example~\ref{ex:u19-lecture}, $\vb{q}_1\T\vb{a}=3/\sqrt3$,
$\vb{q}_1\T\vb{b}=3/\sqrt3$ and $\vb{q}_2\T\vb{b}=2/\sqrt2$, so
\[
  \begin{bmatrix}1&1\\1&0\\1&2\end{bmatrix}
  =\begin{bmatrix}1/\sqrt3&0\\1/\sqrt3&-1/\sqrt2\\1/\sqrt3&1/\sqrt2\end{bmatrix}
   \begin{bmatrix}\sqrt3&\sqrt3\\0&\sqrt2\end{bmatrix}.
\]

\begin{example}[A lucky \texorpdfstring{$3\times 3$}{3 by 3} with no square roots]\label{ex:u19-recitation}
Orthonormalize the columns $\vb{a}=(1,0,0)$, $\vb{b}=(2,0,3)$, $\vb{c}=(4,5,6)$ of
$A=\left[\begin{smallmatrix}1&2&4\\0&0&5\\0&3&6\end{smallmatrix}\right]$, working directly
with unit vectors.
\begin{itemize}
  \item $\vb{a}$ already has length $1$: $\vb{q}_1=(1,0,0)$.
  \item $\vb{b}-(\vb{b}\T\vb{q}_1)\vb{q}_1=(2,0,3)-2(1,0,0)=(0,0,3)$, of length $3$, so
    $\vb{q}_2=(0,0,1)$.
  \item $\vb{c}-(\vb{c}\T\vb{q}_1)\vb{q}_1-(\vb{c}\T\vb{q}_2)\vb{q}_2=(4,5,6)-4(1,0,0)-6(0,0,1)=(0,5,0)$,
    of length $5$, so $\vb{q}_3=(0,1,0)$.
\end{itemize}
The numbers subtracted and the lengths are exactly the entries of $R$:
\[
  \underbrace{\begin{bmatrix}1&2&4\\0&0&5\\0&3&6\end{bmatrix}}_{A}
  =\underbrace{\begin{bmatrix}1&0&0\\0&0&1\\0&1&0\end{bmatrix}}_{Q}
   \underbrace{\begin{bmatrix}1&2&4\\0&3&6\\0&0&5\end{bmatrix}}_{R},
\]
which reads column by column as $\vb{a}=1\vb{q}_1$, $\vb{b}=2\vb{q}_1+3\vb{q}_2$,
$\vb{c}=4\vb{q}_1+6\vb{q}_2+5\vb{q}_3$. Here $Q$ happens to be a permutation matrix, and $R$
is $A$ with rows~2 and~3 exchanged.
\end{example}

\paragraph{Least squares with \texorpdfstring{$A=QR$}{A = QR}.}
Substituting $A=QR$ into the normal equations,
$R\T Q\T QR\hat{\vb{x}}=R\T Q\T\vb{b}$ becomes $R\T R\hat{\vb{x}}=R\T Q\T\vb{b}$, and since $R$
is invertible,
\[
  R\hat{\vb{x}}=Q\T\vb{b},
\]
a triangular system solved by back substitution. This is how least squares is computed in
practice.

\begin{example}[The best-fit line again, via \texorpdfstring{$QR$}{QR}]\label{ex:u19-short}
For the data $(0,6),(1,0),(2,0)$ of Section~\ref{sec:u18}, $A$ has columns
$\vb{a}_1=(1,1,1)$ and $\vb{a}_2=(0,1,2)$, and $\vb{b}=(6,0,0)$. Gram--Schmidt:
$\vb{A}_2=(0,1,2)-\frac33(1,1,1)=(-1,0,1)$, so
\[
  \vb{q}_1=\tfrac{1}{\sqrt3}(1,1,1),\quad \vb{q}_2=\tfrac{1}{\sqrt2}(-1,0,1),\qquad
  R=\begin{bmatrix}\sqrt3&\sqrt3\\0&\sqrt2\end{bmatrix},\qquad
  Q\T\vb{b}=\begin{bmatrix}6/\sqrt3\\-6/\sqrt2\end{bmatrix}=\begin{bmatrix}2\sqrt3\\-3\sqrt2\end{bmatrix}.
\]
Back substitution in $R\hat{\vb{x}}=Q\T\vb{b}$: $\sqrt2\,D=-3\sqrt2$ gives $D=-3$, then
$\sqrt3\,C+\sqrt3\,D=2\sqrt3$ gives $C=5$, the same line $b=5-3t$. The projection is a sum of
two independent one-dimensional projections:
$\vb{p}=\vb{q}_1(\vb{q}_1\T\vb{b})+\vb{q}_2(\vb{q}_2\T\vb{b})=2(1,1,1)-3(-1,0,1)=(5,2,-1)$.
\end{example}

\subsection{Exercises}

\begin{exercise}\label{ex:u19-indep}
Orthonormal vectors are automatically independent. Give a matrix proof: if $Q$ has
orthonormal columns, show that $Q\vb{x}=\vb{0}$ implies $\vb{x}=\vb{0}$. ($Q$ may be
rectangular, so use $Q\T$, not $Q^{-1}$.)
\end{exercise}

\begin{solution}
Multiply $Q\vb{x}=\vb{0}$ by $Q\T$: $Q\T Q\vb{x}=Q\T\vb{0}$, i.e.\ $I\vb{x}=\vb{0}$, so
$\vb{x}=\vb{0}$. Thus $\Nul{Q}=\{\vb{0}\}$ and no nontrivial combination of the columns
vanishes.
\end{solution}

\begin{exercise}\label{ex:u19-perpd}
Apply Gram--Schmidt to $\vb{a}=(1,-1,0,0)$, $\vb{b}=(0,1,-1,0)$, $\vb{c}=(0,0,1,-1)$ to get
orthogonal vectors $\vb{A},\vb{B},\vb{C}$. Show that $\{\vb{A},\vb{B},\vb{C}\}$ and
$\{\vb{a},\vb{b},\vb{c}\}$ are both bases for the space of vectors perpendicular to
$\vb{d}=(1,1,1,1)$.
\end{exercise}

\begin{solution}
$\vb{A}=\vb{a}=(1,-1,0,0)$. Then $\vb{A}\T\vb{b}=-1$, $\vb{A}\T\vb{A}=2$, and
$\vb{B}=\vb{b}+\frac12\vb{A}=(\frac12,\frac12,-1,0)$. Next $\vb{A}\T\vb{c}=0$,
$\vb{B}\T\vb{c}=-1$, $\vb{B}\T\vb{B}=\frac32$, so
$\vb{C}=\vb{c}+\frac23\vb{B}=(\frac13,\frac13,\frac13,-1)$. The three are mutually orthogonal
(check the three dot products) and each has components summing to zero, so each is
perpendicular to $\vb{d}$; so are $\vb{a},\vb{b},\vb{c}$. The space $\vb{d}^{\perp}$ is the
nullspace of the $1\times 4$ matrix $\vb{d}\T$, of dimension $4-1=3$. Three independent
vectors in a $3$-dimensional space form a basis. $\vb{A},\vb{B},\vb{C}$ are independent
because they are orthogonal and nonzero (normalize them and use Exercise~\ref{ex:u19-indep}).
For $\vb{a},\vb{b},\vb{c}$: the combination
$x_1\vb{a}+x_2\vb{b}+x_3\vb{c}=(x_1,\,x_2-x_1,\,x_3-x_2,\,-x_3)$ is zero only when
$x_1=x_2=x_3=0$.
\end{solution}

\begin{exercise}\label{ex:u19-qr}
Find the $QR$ factorization of $A=\begin{bmatrix}1&1\\2&0\\2&1\end{bmatrix}$.
\end{exercise}

\begin{solution}
$\|\vb{a}_1\|=3$, so $\vb{q}_1=\frac13(1,2,2)$ and $r_{11}=3$. Next
$r_{12}=\vb{q}_1\T\vb{a}_2=\frac{1+0+2}{3}=1$, and
$\vb{A}_2=\vb{a}_2-\vb{q}_1=(1,0,1)-\frac13(1,2,2)=\frac13(2,-2,1)$, which already has length
$1$; so $\vb{q}_2=\frac13(2,-2,1)$ and $r_{22}=1$. Hence
\[
  A=\frac13\begin{bmatrix}1&2\\2&-2\\2&1\end{bmatrix}\begin{bmatrix}3&1\\0&1\end{bmatrix},
\]
and multiplying out reproduces the columns $(1,2,2)$ and $(1,0,1)$.
\end{solution}

\begin{exercise}\label{ex:u19-length}
Show that multiplication by a matrix with orthonormal columns preserves lengths and dot
products: $\|Q\vb{x}\|=\|\vb{x}\|$ and $(Q\vb{x})\T(Q\vb{y})=\vb{x}\T\vb{y}$. What does this
say about the rotation matrices of Example~\ref{ex:u19-gallery}?
\end{exercise}

\begin{solution}
$(Q\vb{x})\T(Q\vb{y})=\vb{x}\T Q\T Q\vb{y}=\vb{x}\T\vb{y}$; taking $\vb{y}=\vb{x}$ gives
$\|Q\vb{x}\|^2=\|\vb{x}\|^2$. Rotations therefore move vectors without stretching them and
preserve the angles between them, as a rigid rotation of the plane should.
\end{solution}
